\chapter{Коммутирующий квадрат семейного и наблюдаемого чтений}

Обратный аудит ранних постановок выявил повторяющуюся подмену задачи. Том I
говорил не об одной координатной алгебре, которая обязана породить все поля,
а о нескольких чтениях одного источника. Гейт условных ожиданий Версии III
уже предлагал правильный механизм, но остановился, поскольку происхождение
двух модулей было неизвестно. Теперь этот пробел закрыт:
\begin{equation}
 H_{\mathrm{par}}
 =H_{\mathrm{fam}}\otimes H_{\mathrm{obs}},
 \qquad
 \dim H_{\mathrm{fam}}=20,
 \qquad
 \dim H_{\mathrm{obs}}=15,
 \qquad
 \dim H_{\mathrm{par}}=300.
 \label{eq:v5-commuting-square-factorization}
\end{equation}

Полная матричная алгебра $M_{300}(\mathbb C)$ здесь является не
координатной алгеброй пространства, а окружающей алгеброй общего носителя,
в которой расположены два взаимно коммутирующих чтения.

\section{Две факторные алгебры}

Определим
\begin{equation}
 \mathcal M_{\mathrm{fam}}
 =B(H_{\mathrm{fam}})\otimes I_{15},
 \qquad
 \mathcal M_{\mathrm{obs}}
 =I_{20}\otimes B(H_{\mathrm{obs}}).
\end{equation}
Тогда
\begin{equation}
 [\mathcal M_{\mathrm{fam}},\mathcal M_{\mathrm{obs}}]=0,
 \qquad
 \mathcal M_{\mathrm{fam}}\cap\mathcal M_{\mathrm{obs}}
 =\mathbb C I_{300}.
\end{equation}
Нормированный след факторизуется:
\begin{equation}
 \tau_{300}=\tau_{20}\otimes\tau_{15}.
\end{equation}
Он задаёт канонические условные ожидания
\begin{equation}
 \begin{aligned}
 E_{\mathrm{fam}}&=\operatorname{id}\otimes\tau_{15},\\
 E_{\mathrm{obs}}&=\tau_{20}\otimes\operatorname{id}.
 \end{aligned}
 \label{eq:v5-factor-expectations}
\end{equation}
Оба ожидания унитальны, положительны, сохраняют след и идемпотентны. Они
образуют коммутирующий квадрат:
\begin{equation}
 E_{\mathrm{fam}}E_{\mathrm{obs}}
 =E_{\mathrm{obs}}E_{\mathrm{fam}}
 =\tau_{300}(\,\cdot\,)I_{300}.
 \label{eq:v5-commuting-square}
\end{equation}
Максимальный численный дефект этих тождеств меньше $10^{-14}$.

Именно условные ожидания и коммутирующие квадраты являются стандартным
операторно-алгебраическим способом хранить несколько согласованных
подалгебр одного источника. В теории достаточности Пеца сохраняющее
состояние ожидание также является точной картой обратимого грубого чтения
для выбранного семейства состояний. Но конкретная факторизация
$20\times15$ и её физическая интерпретация являются результатом проекта.

\section{Один след без относительного веса}

Пусть семейная кривизна имеет вид
\begin{equation}
 \mathcal F_{\mathrm{fam}}=[d,d^\dagger]-h.
\end{equation}
Она бесследова:
\begin{equation}
 \tau_{20}(\mathcal F_{\mathrm{fam}})=0,
\end{equation}
поскольку след коммутатора равен нулю, а высота имеет равные крайние ранги.
Для произвольной наблюдаемой кривизны $\mathcal F_{\mathrm{obs}}$ поэтому
\begin{equation}
 \tau_{300}\!\left[
 (\mathcal F_{\mathrm{fam}}\otimes I)
 (I\otimes\mathcal F_{\mathrm{obs}})
 \right]=0.
\end{equation}
Отсюда следует точное разложение
\begin{equation}
 \boxed{
 \tau_{300}(\mathcal F_{\mathrm{tot}}^2)
 =\tau_{20}(\mathcal F_{\mathrm{fam}}^2)
 +\tau_{15}(\mathcal F_{\mathrm{obs}}^2),
 }
 \label{eq:v5-one-trace-pythagoras}
\end{equation}
где
\begin{equation}
 \mathcal F_{\mathrm{tot}}
 =\mathcal F_{\mathrm{fam}}\otimes I
 +I\otimes\mathcal F_{\mathrm{obs}}.
\end{equation}

Таким образом, два сектора не требуют независимого относительного
коэффициента. Их нормы являются ортогональными компонентами одного
гильбертово-шмидтовского квадрата. Старое препятствие гейта условных
ожиданий --- вручную объявленное удвоение модулей --- устранено точной
формулой размерности~\eqref{eq:v5-commuting-square-factorization}.

\section{Какие семейные поля действительно допустимы}

Предыдущий гейт рассматривал общую прямоугольную матрицу
$X\in M_{3\times4}$ и поэтому обнаружил три вращательных нулевых направления.
Но аффинный носитель уже имеет фиксированные действия $S_4$ и $A_4$.

Для перестановочного представления $U_g$ на $\mathbb C^4$ и стандартного
триплета $R_g=VU_gV^T$ пространство сплетающих операторов удовлетворяет
\begin{equation}
 \dim\operatorname{Hom}_{S_4}(\mathbb C^4,\mathbb C^3)=1.
\end{equation}
То же верно после ограничения на $A_4$, а
\begin{equation}
 \dim\operatorname{End}_{A_4}(\mathbb C^3)=1.
\end{equation}
По лемме Шура допустимые стрелки имеют только вид
\begin{equation}
 \boxed{
 X=\rho V,
 \qquad
 Y=\Phi I_3.
 }
 \label{eq:v5-equivariant-field-content}
\end{equation}

Следовательно, непрерывная ориентация триплета не является скалярным полем.
Глобальное вращение одновременно меняет базис представления и матрицу
$V$; это эквивалентное описание одного бимодуля. Вращение только $X$ при
фиксированном $A_4$-действии выходит из пространства допустимых
сплетающих операторов.

\begin{theorem}[Отсутствие ложных семейных мод]
В $S_4/A_4$-эквивариантном классе семейных стрелок три вращательных
направления общего $M_{3\times4}$ не являются полями. Поэтому их не нужно
ни калибровать, ни динамически поднимать тетраэдрическим потенциалом.
\end{theorem}

Это исправляет сделанный нами на предыдущем шаге слишком широкий выбор
пространства флуктуаций. Старый тетраэдрический тензор остаётся важен для
голономии и дискретных осей, но не требуется только ради устранения
искусственно введённых $SO(3)$-нулей.

\section{Сокращённый функционал}

После подстановки~\eqref{eq:v5-equivariant-field-content} три блока
кривизны дают
\begin{equation}
 V(\rho,\Phi)
 =(1-\rho^2)^2
 +(\rho^2-|\Phi|^2)^2
 +( |\Phi|^2-1)^2.
 \label{eq:v5-equivariant-radial-potential}
\end{equation}
При нулевом импульсе минимум
\begin{equation}
 \rho^2=1,
 \qquad |\Phi|^2=1
\end{equation}
имеет по переменным $(\rho,\Re\Phi,\Im\Phi)$ спектр гессиана
\begin{equation}
 \{24,8,0\}.
\end{equation}
Единственный нуль --- фаза $\Phi$, а не семейное вращение.

После добавления единичного дефектного импульса $|\Phi|^2$ получаем тот же
точный минимум, что и в полном прежнем вычислении:
\begin{equation}
 \rho^2=\frac56,
 \qquad
 |\Phi|^2=\frac23,
 \qquad
 V_\star=\frac56.
\end{equation}
Два ненулевых собственных значения равны
\begin{equation}
 12-\frac{4\sqrt{21}}3,
 \qquad
 12+\frac{4\sqrt{21}}3
\end{equation}
и строго положительны.

\section{Что именно найдено}

Новая архитектура не требует, чтобы семейная кривизна была одноформой
алгебры Стандартной модели. Это разные чтения одного оператора на общем
носителе. Их связь задаётся:
\begin{enumerate}
 \item общим состоянием и следом;
 \item коммутирующим квадратом условных ожиданий;
 \item общей кратностью наблюдаемого пакета;
 \item последующими картами чтения в юкавские и дефектные наблюдаемые.
\end{enumerate}

Тем самым снимается ложное требование одной координатной алгебры. Но полное
физическое замыкание ещё не достигнуто. В текущем гейте полная
$B(H_{\mathrm{obs}})$ использовалась как факторный контроль. Следует
заменить её точным представлением
$\mathbb C\oplus\mathbb H\oplus M_3(\mathbb C)$, затем проверить, являются
ли уже найденные сохраняющая состояние проекция юкавского сектора и
фазово-дефектное чтение совместимыми углами того же коммутирующего квадрата.

\begin{protocol}
\begin{enumerate}
 \item происхождение двух модулей старого гейта условных ожиданий:
 \textbf{получено};
 \item общий носитель размерности $300$:
 \textbf{получен};
 \item два коммутирующих чтения и канонические условные ожидания:
 \textbf{получены};
 \item один след без относительного веса секторов:
 \textbf{получен};
 \item $S_4/A_4$-эквивариантное пространство стрелок:
 \textbf{одномерно для каждой стрелки};
 \item три вращательных нуля как физические поля:
 \textbf{исключены};
 \item устойчивый радиальный вакуум при единичном импульсе:
 \textbf{сохранён};
 \item точное наблюдаемое представление Стандартной модели в квадрате:
 \textbf{ещё не подставлено};
 \item карта к юкавским и фазовым наблюдаемым:
 \textbf{открыта};
 \item физическое замыкание:
 \textbf{не достигнуто}.
\end{enumerate}
\end{protocol}

Машинная проверка и сертификат сохранены в
\path{s2t_v5_commuting_square_readout_gate.py} и
\path{s2t_v5_commuting_square_readout_gate_results.json}.